\section{Overview and principal findings}
\label{sec:verdict}

\begin{verdictbox}[Principal finding]
The proposal is feasible for a single developer working with
LLM agents, subject to one architectural correction. A machine-checked
proof layer over \dotnetrs{}'s \safetyc{} corpus is tractable because
the corpus factors into approximately nine invariant families,
40--60 proof sites of genuine difficulty, and a small number of
protocol theorems --- not 620 independent obligations. The correction
concerns the checker: the dependently-typed language should be a
domain-fitted \emph{surface} --- an embedded, tactic-oriented
proof-script DSL in the lineage of \code{dotnetdll}'s \code{asm!}
macro, whose premises are ghost tokens checked by \code{rustc} and
whose scripts elaborate into Lean 4 --- rather than a calculus with its
own kernel and solver. The fallback position named in the original
proposal (``develop a model environment in Lean or Rocq and compile
our proof language to that'') is, on the evidence assembled here, the
correct primary design.
\end{verdictbox}

Six findings carry the argument:

\begin{enumerate}[leftmargin=1.6em]
\item \textbf{The premise is confirmed empirically in this
  repository.} Two soundness defects entered the tree in the five
  commits preceding this study --- a feature-gated no-op cited as an
  alignment witness, and a \code{Sync} implementation justified by a
  false claim --- and both instantiate a single failure mode:
  \emph{prose asserting a witness that no code establishes}. Under
  the ghost-token design the first is not expressible --- citing a
  \code{cfg}-excluded guard is a compile error --- and the proof layer
  addresses the class (\cref{sec:system-defects},
  \cref{sec:dsl-example-fail}). Both defects, and the rest of Phase 0,
  have since been closed by ordinary review (\code{208a6c8b}); the
  premise is unweakened by that fix landing without a proof layer ---
  it argues for the class of catch this proposal targets, not for any
  particular open instance of it.
\item \textbf{Each component has substantial precedent; the
  composition does not.} Machine-checked proofs of real \code{unsafe}
  Rust exist (RefinedRust, VeriFast, Verus); machine-readable safety
  annotations are an emerging Rust convention (RFC 3842,
  \code{tag-std}); verified collectors and verified bytecode machines
  exist (Hawblitzel--Petrank, JinjaThreads). What exists nowhere is a
  connection between such proofs and a managed-runtime specification
  model: \emph{no mechanized formalization of ECMA-335 exists}. That
  gap accounts for roughly half the projected cost and all of the
  novelty; an executable Lean model of the VES core would be a
  publishable artifact independent of the proof campaign
  (\cref{sec:prior-art}).
\item \textbf{The verification problem is two-level by nature.} The
  difficult \safetyc{} comments appeal simultaneously to Rust-side
  facts (validity, aliasing, provenance under Tree Borrows) and
  VES-side facts (stop-the-world protocol state, layout faithfulness,
  evaluation-stack typing, GC reachability). The architecture must
  therefore pair a Rust memory-model fragment with an ECMA-335
  abstract machine and a representation relation between them --- the
  structure of a compiler-correctness proof, applied in the reverse
  direction (\cref{sec:theory}, \cref{sec:design-model}).
\item \textbf{Lean 4 is the appropriate primary backend; Rocq is the
  appropriate escalation path.} Lean is stronger on the model half
  (executable specifications, differential testing against the
  existing 486-fixture oracle, elaborator metaprogramming for the DSL,
  the largest automation and agent-tooling ecosystem); Rocq holds the
  foundational Rust-semantics results (Iris, RustBelt, Tree Borrows,
  RefinedRust). The recommended design states Rust-substrate
  obligations against a small, Miri-falsifiable axiomatization in
  Lean and retains RefinedRust-in-Rocq as an upgrade path for a small
  set of critical modules (\cref{sec:design-embed}).
\item \textbf{Compartmentalization determines survivability.} The
  verifier decomposes into five small Rust crates and five Lean
  packages, each separately versioned, publishable, and abandonable
  (\cref{sec:design-compart}); the only artifacts crossing the
  boundary to \dotnetrs{} are two small build dependencies and two
  content-hashed manifests. The VM's build and refactor velocity
  therefore never couple to Lean, and a stalled proof campaign
  freezes sites at \code{trusted} rather than breaking anything.
\item \textbf{Staging dominates ambition.} The roadmap
  (\cref{sec:roadmap}) is cut so that every phase delivers an
  artifact the VM uses in CI within that phase: hygiene fixes and a
  blocking UB gate; the embedded token and script layer with
  \code{rustc}-checked premises; the executable VES model; the STW
  protocol theorem (the principal technical risk,
  \cref{sec:risks-concurrency}); and finally the site campaign under
  a trusted-premise ratchet. \Cref{sec:decomposition} refines the
  phases into an ordered sequence of supervised agent-refactor work
  packages sized by checklist complexity and verification-signal
  quality rather than by calendar time.
\end{enumerate}
